PASTE STUFF HERE

THIS IS FOR THE DISCHARGIN CASE OF STEP SIGNAL

this is for the question 2 first condition , where t<RC

\subsection*{Question 2: Short pulse ($T_p \ll RC$)}
\subsubsection*{Hand Calculation}
\[
T_p \ll 0.82\,ms
\]
Peak capacitor voltage:
\begin{equation}
V_{max} = V_s\left(1 - e^{-T_p/RC}\right)
\end{equation}
For $T_p \ll RC$:
\[
V_{max} \approx \frac{T_p}{RC}V_s
\]

2nd condition 

\subsection*{Question 2: Pulse width $T_p \approx RC$}
\subsubsection*{Hand Calculation}
\[
T_p \approx 0.82\,ms
\]
\[
V_C = V_s(1 - e^{-1}) = 0.632V_s
\]

3rd condition 

\subsection*{Question 2: Long pulse ($T_p \gg RC$)}
\subsubsection*{Hand Calculation}
\[
T_p \gg 0.82\,ms
\]
\[
V_C \approx V_s
\]

Capacitor charges almost completely during ON time and discharges exponentially during OFF time.


THIS IS FOR THE 3RD PART 
\section{Part C: Square-Wave Input (Transient and Steady-State)}
\subsection{Case 1: $RC \gg T$}

\[
T \ll RC
\]

\[
T \ll 0.82\,ms
\]


\subsubsection*{Hand Calculations}

For capacitor charging:

\begin{equation}
V_C(t) = V_i + (V_f - V_i)\left(1 - e^{-t/RC}\right)
\end{equation}

Since $t \le \frac{T}{2} \ll RC$,

\[
\frac{t}{RC} \ll 1
\]

Using approximation:

\[
e^{-t/RC} \approx 1 - \frac{t}{RC}
\]

\[
V_C(t) \approx V_i + (V_f - V_i)\frac{t}{RC}
\]

\subsubsection*{Inference}

The capacitor voltage varies almost linearly with time.  
Hence, the output waveform is approximately triangular.

\[
\boxed{\text{RC circuit behaves as an integrator}}
\]

THIS IS THE SECOND PART OF 3RD QUESTION
\subsection{Case 2: $RC \approx T$}

\[
T \approx RC = 0.82\,ms
\]

\subsubsection*{Hand Calculations}

During the charging interval:

\begin{equation}
V_C(t) = V_s\left(1 - e^{-t/RC}\right)
\end{equation}

At the end of half cycle:

\[
t = \frac{T}{2} \approx \frac{RC}{2}
\]

\[
V_C = V_s\left(1 - e^{-0.5}\right)
\]

\[
V_C \approx 0.393V_s
\]

During discharge:

\begin{equation}
V_C(t) = V_0 e^{-t/RC}
\end{equation}

At $t = \frac{T}{2}$:

\[
V_C = V_0 e^{-0.5} \approx 0.607V_0
\]

\subsubsection*{Inference}

The capacitor partially charges and partially discharges in each cycle.  
Both transient and steady-state components are present.

\[
\boxed{\text{Maximum transient behavior is observed}}
\]

THIS IS CASE THREE
\subsection{Case 3: $RC \ll T$}

\[
T \gg RC
\]

\[
T \gg 0.82\,ms
\]

\subsubsection*{Hand Calculations}

During ON time:

\[
t = \frac{T}{2} \gg RC
\]

\[
e^{-t/RC} \approx 0
\]

\[
V_C \approx V_s
\]

During OFF time:

\begin{equation}
V_C(t) = V_s e^{-t/RC}
\end{equation}

For $t \gg RC$:

\[
V_C \approx 0
\]

\subsubsection*{Inference}

The capacitor fully charges and discharges within each cycle.

\[
\boxed{\text{Output closely follows the input waveform}}
\]


VERIFICATION USING TIME CONSTANT
\subsubsection{Verification Using $2\tau$}

The theoretical time constant of the RC circuit is:

\[
\tau = RC = 820 \times 1\times10^{-6} = 0.82\,ms
\]

Hence,

\[
2\tau = 2 \times 0.82 = 1.64\,ms
\]

For a step input, the capacitor charging equation is:

\[
V_C(t) = V_s\left(1 - e^{-t/\tau}\right)
\]

At \(t = 2\tau\):

\[
V_C(2\tau) = V_s\left(1 - e^{-2}\right)
\]

\[
V_C(2\tau) = V_s(1 - 0.135)
\]

\[
\boxed{V_C(2\tau) \approx 0.865\,V_s}
\]

From the experimental waveform, the capacitor voltage at approximately
\(t = 1.64\,ms\) was observed to be close to \(86.5\%\) of the final value.

Thus, the experimental result agrees well with the theoretical prediction, verifying the step response of the RC circuit.

